\section{Bridge}
\label{sec:pat-bridge}

\problemstatement{Decouple an abstraction from its implementation so the
two can vary independently, avoiding a combinatorial explosion of
subclasses when a type varies along two separate dimensions.}

\begin{patternbox}[When You Reach for It]
The trigger is \emph{``two independent dimensions of variation''} that,
combined by inheritance, would create $m\times n$ subclasses. Shapes
(circle, square) times rendering APIs (OpenGL, DirectX) would need
\texttt{OpenGLCircle}, \texttt{DirectXCircle}, \texttt{OpenGLSquare}\ldots
Bridge turns that product into a sum.
\end{patternbox}

\subsection*{Understanding the pattern}

Split the two dimensions into two class hierarchies and connect them with
a reference (the ``bridge''). The \emph{abstraction} hierarchy (shapes)
holds a pointer to an \emph{implementor} hierarchy (renderers) and
delegates the primitive operations. Now you add shapes and renderers
independently: $m+n$ classes instead of $m\times n$.

\begin{ideabox}[Prefer Composition Across the Two Axes]
When a class varies along two axes, do not multiply them with inheritance.
Give one axis (the abstraction) a member pointing at the other (the
implementation). This is ``favour composition over inheritance'' applied
to prevent a subclass explosion.
\end{ideabox}

\subsection*{C++ implementation}

\begin{lstlisting}
#include <bits/stdc++.h>
using namespace std;

// Implementor axis: how to render.
struct Renderer {
    virtual void renderCircle(float r) = 0;
    virtual ~Renderer() = default;
};
struct OpenGLRenderer : Renderer {
    void renderCircle(float r) override { cout << "OpenGL circle r=" << r << "\n"; }
};
struct DirectXRenderer : Renderer {
    void renderCircle(float r) override { cout << "DirectX circle r=" << r << "\n"; }
};

// Abstraction axis: what to draw. Holds a bridge to the implementor.
class Shape {
protected:
    Renderer& renderer;                        // the bridge
public:
    Shape(Renderer& r) : renderer(r) {}
    virtual void draw() = 0;
    virtual ~Shape() = default;
};
class Circle : public Shape {
    float radius;
public:
    Circle(Renderer& r, float rad) : Shape(r), radius(rad) {}
    void draw() override { renderer.renderCircle(radius); }
};

int main() {
    OpenGLRenderer gl; DirectXRenderer dx;
    Circle c1(gl, 5), c2(dx, 8);               // mix axes freely
    c1.draw();
    c2.draw();
    return 0;
}
\end{lstlisting}

\begin{complexitybox}[Trade-offs]
\textbf{Pros.} Abstraction and implementation evolve independently; avoids
class explosion; swappable implementations at run time. \textbf{Cons.}
More indirection and up-front structure; only pays off when \emph{both}
dimensions genuinely vary. For a single axis it is over-engineering.
\end{complexitybox}

\begin{takeawaybox}
Bridge separates ``what'' from ``how'' into two hierarchies joined by a
reference, turning an $m\times n$ subclass explosion into $m+n$. The
signal is two orthogonal dimensions of change. It looks structurally like
Strategy but its intent is structural decoupling, not swapping an
algorithm.
\end{takeawaybox}
